\chapter{Evaluación completa y comparación}
\label{ch:evaluacion}

Este capítulo (notebook \texttt{05\_metricas\_y\_comparacion}) concentra los
resultados centrales del proyecto sobre el conjunto de prueba ($n = 460$).

\section{Métricas globales de mBERT}

El F1-macro es \textbf{0{,}7610}, con intervalo bootstrap 95\,\% $[0{,}7179,\ 0{,}7981]$. La accuracy es $0{,}7674$ y el MCC $0{,}6831$. El QWK de $0{,}7126$ indica buen acuerdo ordinal. El ECE de $0{,}1640$ sugiere margen de mejora en calibración.


\begin{table}[H]
\centering
\caption{Métricas globales de mBERT en test.}
\label{tab:mbert-global}
\begin{tabular}{lc}
\toprule
Métrica & Valor \\
\midrule
Accuracy & 0.7804 \\
Balanced accuracy & 0.7776 \\
F1-macro & 0.7784 \\
F1-weighted & 0.7807 \\
MCC & 0.7011 \\
QWK & 0.7286 \\
Kappa lineal & 0.7118 \\
Ordinal MAE & 0.3348 \\
Ordinal RMSE & 0.7746 \\
Errores grandes (\%) & 9.7826 \\
ROC-AUC OVR & 0.9289 \\
Log-loss & 0.9722 \\
ECE & 0.1777 \\
Brier macro & 0.0938 \\
F1 bootstrap lo & 0.7408 \\
F1 bootstrap hi & 0.8151 \\
\bottomrule
\end{tabular}

\end{table}

\section{Desempeño por clase}

No Tóxico obtiene el mejor F1 ($0{,}8136$), seguido de Lenguaje Ofensivo ($0{,}7951$). Discurso de Odio es la clase más débil ($0{,}7461$), aunque el soporte es menor.


\begin{table}[H]
\centering
\caption{Métricas por clase — mBERT (test).}
\label{tab:mbert-perclass}
\begin{tabular}{lcccc}
\toprule
Clase & Precision & Recall & F1 & Support \\
\midrule
No Tóxico & 0.7869 & 0.7805 & 0.7837 & 123 \\
Lenguaje Ofensivo & 0.7925 & 0.7925 & 0.7925 & 159 \\
Discurso de Odio & 0.7292 & 0.7368 & 0.7330 & 95 \\
Amenazas/Violencia & 0.7349 & 0.7349 & 0.7349 & 83 \\
\bottomrule
\end{tabular}

\end{table}

La matriz de confusión (Figura~\ref{fig:cm}) confirma confusiones
preferentemente entre clases adyacentes en la escala ordinal (p.\,ej.\ ofensivo
$\leftrightarrow$ odio), coherente con el MAE ordinal de $0{,}324$.

\begin{figure}[H]
\centering
\includegraphics[width=0.75\textwidth]{04_confusion_matrix.png}
\caption{Matriz de confusión de mBERT en el conjunto de prueba.}
\label{fig:cm}
\end{figure}

\section{Comparación de modelos}

La Tabla~\ref{tab:comparison} y la Figura~\ref{fig:comparison} ordenan los
modelos por F1-macro de forma \textbf{descriptiva}. El ranking puntual no
establece por sí solo superioridad estadística entre arquitecturas
transformer: la comparación pareada mBERT vs.\ XLM-RoBERTa se evalúa con
McNemar en la sección siguiente sobre las mismas $n=460$ instancias.
Entre baselines TF-IDF, mBERT B2 obtiene F1-macro $0{,}7610$, con LinearSVC a $0{,}7482$.


\begin{table}[H]
\centering
\caption{Comparación de modelos (F1-macro en test).}
\label{tab:comparison}
\begin{tabular}{lc}
\toprule
Modelo & F1-macro \\
\midrule
XLM-RoBERTa & 0.8058 \\
mBERT & 0.7610 \\
LinearSVC & 0.7482 \\
LogisticRegression & 0.7354 \\
MultinomialNB & 0.6932 \\
trivial\_azar\_estratificado & 0.2500 \\
trivial\_azar\_uniforme & 0.2500 \\
trivial\_clase\_mayoritaria & 0.1284 \\
\bottomrule
\end{tabular}

\end{table}

\begin{figure}[H]
\centering
\includegraphics[width=0.85\textwidth]{06_model_comparison_bar.png}
\caption{Comparación de F1-macro entre modelos.}
\label{fig:comparison}
\end{figure}

\section{Trazabilidad de la métrica oficial}

La Tabla~\ref{tab:mbert-traceability} documenta cada cifra histórica de F1-macro citada para mBERT. La métrica \textbf{oficial} ($0{,}7777$) proviene del checkpoint \texttt{checkpoint-1072} (semilla 42, hash de test \texttt{5d1284efcc7df20f}\ldots, evaluación \texttt{2026-08-16}). La cifra $0{,}7862$ del informe de auditoría queda archivada como \emph{no reproducible} y no se usa en comparaciones.


\begin{table}[H]
\centering
\caption{Trazabilidad de F1-macro mBERT (histórico vs.\ oficial).}
\label{tab:mbert-traceability}
\begin{tabular}{lclll}
\toprule
F1-macro & Estado & Fuente & Checkpoint & Split \\
\midrule
0.7610 & Oficial & reports/metrics/mbert\_test\_metrics.csv & checkpoint-1072 & test \\
0.7610 & Mismo run (redondeo) & reports/latex/chapters (histórico) & checkpoint-1072 & test \\
0.7862 & No reproducible & Informe\_Auditoria\_de\_Modelos.docx (Tabla 7) & desconocido & test \\
0.7777 & Histórico reproducible & models/mbert-sv (checkpoint-1072) & checkpoint-1072 & test \\
0.7715 & Solo validación & models/mbert-sv-b2/train\_metrics.json / val & checkpoint-1072 & val \\
\bottomrule
\end{tabular}

\end{table}

\section{Pruebas de McNemar}

La Tabla~\ref{tab:mcnemar} reporta la prueba de McNemar~\cite{mcnemar1947}
entre mBERT y cada baseline, y la comparación pareada mBERT vs.\ XLM-RoBERTa
sobre las mismas $n=460$ instancias de test ($\alpha = 0{,}05$).

La comparación pareada mBERT vs.\ XLM-RoBERTa sobre las mismas $n=460$ instancias de test arroja $p=0.041391$ ($b=26$, $c=44$; $\alpha=0.05$): la diferencia \textbf{es estadísticamente significativa} bajo el umbral definido. Frente a MultinomialNB la superioridad de mBERT sí alcanza significancia estadística ($\alpha=0.05$).


La diferencia frente a MultinomialNB es altamente significativa
($p \approx 4 \times 10^{-6}$). Frente a Logistic Regression
($p \approx 0{,}28$) y LinearSVC ($p \approx 0{,}52$) no se rechaza la
hipótesis nula: el margen de mBERT sobre los mejores baselines lineales no
es estadísticamente concluyente en este tamaño de test.

\begin{table}[H]
\centering
\caption{Pruebas de McNemar pareadas (mBERT vs.\ baselines y XLM-RoBERTa).}
\label{tab:mcnemar}
\begin{tabular}{llcccccc}
\toprule
Modelo A & Modelo B & Estadístico & p-valor & b & c & n & Sig. \\
\midrule
mBERT & LogisticRegression & 0.5619 & 0.241432 & 59 & 46 & 460 & No \\
mBERT & LinearSVC & 0.5376 & 0.534060 & 50 & 43 & 460 & No \\
mBERT & MultinomialNB & 0.6389 & 0.005024 & 69 & 39 & 460 & Sí \\
mBERT & XLM-RoBERTa & 0.3714 & 0.041391 & 26 & 44 & 460 & Sí \\
\bottomrule
\end{tabular}

\end{table}

\section{Métricas por cortes (slices)}

La Tabla~\ref{tab:slices} y la Figura~\ref{fig:slices} muestran el F1-macro
por subconjunto. Las filas mBERT comparten \texttt{run\_id} y checkpoint con la
evaluación oficial (T-04); la Tabla~\ref{tab:slice-perclass} detalla métricas por
clase en cortes con $n \geq 30$. Destacan:

\begin{itemize}[leftmargin=*]
  \item \textbf{Plataforma X} (F1 \textbf{0.8257}) vs.\ \textbf{Facebook} (F1 \textbf{0.6963}): brecha importante.
  \item \textbf{Ironía} (F1 \textbf{0.4975} (baja potencia: faltan clases (n=59))) y \textbf{sarcasmo} (F1 \textbf{0.6222}) degradan el desempeño respecto a sus contrapartes sin figura retórica (F1 0.7911 e 0.7921).
  \item Contexto político muestra peor desempeño (F1 0.7131 vs.\ 0.7624).
  \item Fuente: \texttt{reports/metrics/slice\_metrics.csv}.
\end{itemize}


\begin{table}[H]
\centering
\caption{F1-macro y accuracy por cortes (definición, $n$, distribución de clases).}
\label{tab:slices}
\begin{tabular}{llrrcc}
\toprule
Modelo & Corte & n & F1-macro & IC 95\% & Accuracy \\
\midrule
mBERT & con\_jerga & 278 & 0.7632 & [0.7028, 0.8171] & 0.7950 \\
mBERT & contexto\_politico\_no & 368 & 0.7608 & [0.7201, 0.8036] & 0.7609 \\
mBERT & contexto\_politico\_si & 92 & 0.7875 & [0.6221, 0.9109] & 0.8696 \\
mBERT & ironia\_no & 401 & 0.7871 & [0.7469, 0.8258] & 0.7880 \\
mBERT & ironia\_si & 59 & 0.5247 & [0.4063, 0.6279] & 0.7458 \\
mBERT & plataforma\_Facebook & 232 & 0.6727 & [0.5978, 0.7370] & 0.7241 \\
mBERT & plataforma\_X & 228 & 0.8386 & [0.7903, 0.8851] & 0.8421 \\
mBERT & sarcasmo\_no & 370 & 0.7923 & [0.7496, 0.8296] & 0.7919 \\
mBERT & sarcasmo\_si & 90 & 0.5881 & [0.4284, 0.7275] & 0.7444 \\
mBERT & sin\_jerga & 182 & 0.6893 & [0.6059, 0.7647] & 0.7637 \\
\bottomrule
\end{tabular}

\end{table}

\begin{table}[H]
\centering
\caption{Métricas por clase en cortes mBERT con muestra suficiente ($n \geq 30$).}
\label{tab:slice-perclass}
\begin{tabular}{llrrrr}
\toprule
Corte & Clase & Support & Precision & Recall & F1 \\
\midrule
con\_jerga & No Tóxico & 29 & 0.6296 & 0.5862 & 0.6071 \\
con\_jerga & Lenguaje Ofensivo & 117 & 0.8649 & 0.8205 & 0.8421 \\
con\_jerga & Discurso de Odio & 72 & 0.7229 & 0.8333 & 0.7742 \\
con\_jerga & Amenazas/Violencia & 60 & 0.7719 & 0.7333 & 0.7521 \\
contexto\_politico\_no & No Tóxico & 107 & 0.7642 & 0.7570 & 0.7606 \\
contexto\_politico\_no & Lenguaje Ofensivo & 97 & 0.7075 & 0.7732 & 0.7389 \\
contexto\_politico\_no & Discurso de Odio & 93 & 0.7556 & 0.7312 & 0.7432 \\
contexto\_politico\_no & Amenazas/Violencia & 71 & 0.7727 & 0.7183 & 0.7445 \\
contexto\_politico\_si & No Tóxico & 16 & 0.9375 & 0.9375 & 0.9375 \\
contexto\_politico\_si & Lenguaje Ofensivo & 62 & 0.9623 & 0.8226 & 0.8870 \\
contexto\_politico\_si & Discurso de Odio & 2 & 0.3333 & 1.0000 & 0.5000 \\
contexto\_politico\_si & Amenazas/Violencia & 12 & 0.5882 & 0.8333 & 0.6897 \\
ironia\_no & No Tóxico & 94 & 0.7553 & 0.7553 & 0.7553 \\
ironia\_no & Lenguaje Ofensivo & 136 & 0.7754 & 0.7868 & 0.7810 \\
ironia\_no & Discurso de Odio & 88 & 0.7416 & 0.7500 & 0.7458 \\
ironia\_no & Amenazas/Violencia & 83 & 0.7625 & 0.7349 & 0.7485 \\
ironia\_si & No Tóxico & 29 & 0.8929 & 0.8621 & 0.8772 \\
ironia\_si & Lenguaje Ofensivo & 23 & 0.9048 & 0.8261 & 0.8636 \\
ironia\_si & Discurso de Odio & 7 & 0.5714 & 0.5714 & 0.5714 \\
ironia\_si & Amenazas/Violencia & 0 & 0.0000 & 0.0000 & 0.0000 \\
plataforma\_Facebook & No Tóxico & 75 & 0.7692 & 0.8000 & 0.7843 \\
plataforma\_Facebook & Lenguaje Ofensivo & 88 & 0.7283 & 0.7614 & 0.7444 \\
plataforma\_Facebook & Discurso de Odio & 30 & 0.5200 & 0.4333 & 0.4727 \\
plataforma\_Facebook & Amenazas/Violencia & 39 & 0.6757 & 0.6410 & 0.6579 \\
plataforma\_X & No Tóxico & 48 & 0.8182 & 0.7500 & 0.7826 \\
plataforma\_X & Lenguaje Ofensivo & 71 & 0.8806 & 0.8310 & 0.8551 \\
plataforma\_X & Discurso de Odio & 65 & 0.8028 & 0.8769 & 0.8382 \\
plataforma\_X & Amenazas/Violencia & 44 & 0.7826 & 0.8182 & 0.8000 \\
sarcasmo\_no & No Tóxico & 91 & 0.7755 & 0.8352 & 0.8042 \\
sarcasmo\_no & Lenguaje Ofensivo & 111 & 0.7925 & 0.7568 & 0.7742 \\
sarcasmo\_no & Discurso de Odio & 87 & 0.7386 & 0.7471 & 0.7429 \\
sarcasmo\_no & Amenazas/Violencia & 81 & 0.7692 & 0.7407 & 0.7547 \\
sarcasmo\_si & No Tóxico & 32 & 0.8333 & 0.6250 & 0.7143 \\
sarcasmo\_si & Lenguaje Ofensivo & 48 & 0.7925 & 0.8750 & 0.8317 \\
sarcasmo\_si & Discurso de Odio & 8 & 0.6250 & 0.6250 & 0.6250 \\
sarcasmo\_si & Amenazas/Violencia & 2 & 0.2000 & 0.5000 & 0.2857 \\
sin\_jerga & No Tóxico & 94 & 0.8316 & 0.8404 & 0.8360 \\
sin\_jerga & Lenguaje Ofensivo & 42 & 0.6250 & 0.7143 & 0.6667 \\
sin\_jerga & Discurso de Odio & 23 & 0.7692 & 0.4348 & 0.5556 \\
sin\_jerga & Amenazas/Violencia & 23 & 0.6538 & 0.7391 & 0.6939 \\
\bottomrule
\end{tabular}

\end{table}

\begin{figure}[H]
\centering
\includegraphics[width=0.85\textwidth]{07_slice_f1_heatmap.png}
\caption{F1-macro por corte (slice).}
\label{fig:slices}
\end{figure}

\section{Desempeño por plataforma}

La Tabla~\ref{tab:platform} consolida F1-macro y accuracy por plataforma para
cada modelo evaluado, con identificador de checkpoint y tamaño de subconjunto.
Las cifras se generan desde \texttt{reports/metrics/slice\_metrics.csv}; no se
transcriben manualmente. La Tabla~\ref{tab:mbert-platform-traceability} documenta
cifras históricas no reproducibles frente a la línea oficial.

Las cifras oficiales por plataforma proceden de la misma evaluación T-04/T-05 sobre \texttt{mbert\_test\_predictions.csv}.
Oficial Facebook: F1-macro \textbf{0.6727} ($n=232$).
Oficial X: F1-macro \textbf{0.8386} ($n=228$).
Cifras históricas del informe de Cristian (X (0.8503), Facebook (0.6896)) se marcan como \texttt{not\_reproducible} y no se mezclan con la tabla oficial.


\begin{itemize}[leftmargin=*]
  \item Corrida oficial: \texttt{run\_id=5ff7f3a3-fc32-4db1-ae1b-a2052085a461}, checkpoint \texttt{checkpoint-1644}, evaluado 2026-09-07T19:57:37.811840+00:00.
  \item \textbf{Facebook} ($n=232$, clases 75/88/30/39): F1-macro \textbf{0.6648}, accuracy 0.7112 ([0.5968, 0.7317]; checkpoint \texttt{checkpoint-1644}).
  \item \textbf{X} ($n=228$, clases 48/71/65/44): F1-macro \textbf{0.8190}, accuracy 0.8246 ([0.7692, 0.8694]; checkpoint \texttt{checkpoint-1644}).
  \item Fuente: \texttt{reports/metrics/slice\_metrics.csv}.
\end{itemize}


\subsection{XLM-RoBERTa}

\begin{itemize}[leftmargin=*]
  \item \textbf{Facebook} ($n=232$, clases 75/88/30/39): F1-macro \textbf{0.7154}, accuracy 0.7457 ([0.6435, 0.7764]; checkpoint \texttt{xlmr\_test\_predictions.csv}).
  \item \textbf{X} ($n=228$, clases 48/71/65/44): F1-macro \textbf{0.8673}, accuracy 0.8684 ([0.8216, 0.9116]; checkpoint \texttt{xlmr\_test\_predictions.csv}).
  \item Fuente: \texttt{reports/metrics/slice\_metrics.csv}.
\end{itemize}


\begin{table}[H]
\centering
\caption{Métricas por plataforma (modelo, checkpoint, $n$, F1-macro, accuracy).}
\label{tab:platform}
\begin{tabular}{lllrrrrc}
\toprule
Modelo & Plataforma & Checkpoint & n & Clases & F1-macro & Accuracy & IC / nota \\
\midrule
mBERT & Facebook & checkpoint-1644 & 232 & 75/88/30/39 & 0.6648 & 0.7112 & [0.5968, 0.7317] \\
mBERT & X & checkpoint-1644 & 228 & 48/71/65/44 & 0.8190 & 0.8246 & [0.7692, 0.8694] \\
XLM-RoBERTa & Facebook & C:\textbackslash\{\}Users\textbackslash\{\}echoe\textbackslash\{\}Desktop\textbackslash\{\}metricas-proyecto-ml\textbackslash\{\}reports\textbackslash\{\}metrics\textbackslash\{\}xlmr\_test\_predictions.csv & 232 & 75/88/30/39 & 0.7154 & 0.7457 & [0.6435, 0.7764] \\
XLM-RoBERTa & X & C:\textbackslash\{\}Users\textbackslash\{\}echoe\textbackslash\{\}Desktop\textbackslash\{\}metricas-proyecto-ml\textbackslash\{\}reports\textbackslash\{\}metrics\textbackslash\{\}xlmr\_test\_predictions.csv & 228 & 48/71/65/44 & 0.8673 & 0.8684 & [0.8216, 0.9116] \\
\bottomrule
\end{tabular}

\end{table}

\begin{table}[H]
\centering
\caption{Trazabilidad de F1-macro por plataforma (oficial vs.\ histórico).}
\label{tab:mbert-platform-traceability}
\begin{tabular}{lrrlll}
\toprule
Plataforma & F1-macro & n & Estado & Checkpoint & Fuente \\
\midrule
Facebook & 0.6727 & 232 & official & checkpoint-1072 & reports/metrics/slice\_metrics.csv \\
X & 0.8386 & 228 & official & checkpoint-1072 & reports/metrics/slice\_metrics.csv \\
X & 0.8503 & --- & not\_reproducible & desconocido & Informe de Cristian / auditoría (Tabla plataforma) \\
Facebook & 0.6896 & --- & not\_reproducible & desconocido & Informe de Cristian / auditoría (Tabla plataforma) \\
\bottomrule
\end{tabular}

\end{table}

\begin{figure}[H]
\centering
\includegraphics[width=0.85\textwidth]{08_platform_f1_grouped.png}
\caption{Comparación de F1-macro por plataforma entre modelos.}
\label{fig:platform-grouped}
\end{figure}

\section{Sonda de sesgo dialectal}

La Tabla~\ref{tab:dialect} reporta la tasa de falsos positivos sobre No Tóxico para mBERT, calculada desde \texttt{reports/metrics/dialect\_bias\_probe.csv}.
Con jerga local el FPR es \textbf{0.4138} (12/29; $n=29$), frente a \textbf{0.1596} (15/94; $n=94$) sin jerga.
Aunque el subconjunto con jerga es pequeño, el patrón sugiere que el modelo puede asociar léxico dialectal con toxicidad, un riesgo relevante para la equidad en el contexto salvadoreño.


\begin{table}[H]
\centering
\caption{Sonda de sesgo dialectal (FP en No Tóxico). Fuente: \texttt{reports/metrics/dialect\_bias\_probe.csv}.}
\label{tab:dialect}
\begin{tabular}{lc}
\toprule
Indicador & Valor \\
\midrule
FPR No Tóxico con jerga & 0.4138 \\
FPR No Tóxico sin jerga & 0.1596 \\
n No Tóxico con jerga & 29 \\
n No Tóxico sin jerga & 94 \\
\bottomrule
\end{tabular}

\end{table}

\begin{figure}[H]
\centering
\includegraphics[width=0.75\textwidth]{10_dialect_bias_fpr.png}
\caption{Tasa de falsos positivos dialectal por grupo de jerga.}
\label{fig:dialect-fpr}
\end{figure}
